
From the previous result, we notice that BERT model perform much better than the GPT model in the spam email detection task. 
This is because of the following reasons.

1. Fundamental architectural differences between the two models.

We know that BERT Uses Bidirectional Context, meaning it considers both the left and right 
context when encoding a word. This allows BERT to capture rich contextual information about the entire email, making it well-suited for classification tasks where understanding the whole text is important.

GPT-2 is an autoregressive model, meaning it generates text from left to right, only considering past tokens. 
This limits its ability to understand the full context of an email, especially when key spam indicators appear later in the text.

So from the result, we can say that BERT is more suitable to handle classification tasks than GPT-2.

2. Effective Epochs and Training Time

In our experience, even with 20 epochs, BERT can run much faster, in around 10 minutes, while for GPT-2 model, the whole processing time will around 30 minutes, even if we set the per\_device\_train\_batch\_size and gradient\_accumulation\_steps to a relative large number, which is 64 and 2, respectively.

In the GPT-2 experiment, we also found that after 2 epoches, the test loss began to increase, which make sense, because in most of the LLM supervised training, we will only train the input 2-4 round, not 20 in our case, which lead to the overfiting issue in our GPT-2 experiment.

Regarding the limitations of this work, we can point out the following:

1. Limited Dataset

We only use the text data instead of html mail which may contains more multimedia information, such as picture and video, so the model we use may not be able to detect spam emails with multimedia information.

2. Limited Models

We only use the BERT and GPT model to detect spam emails, and there are many other models that can be used to detect spam emails, like DistilBERT \cite{sanh2019distilbert}, RoBERTa \cite{liu2019roberta}, ALBERT \cite{lan2019albert}, which may generate better performance than BERT.

About the future work and the limitations mentioned above, we need use more advanced models to detect spam emails, to identify those spam emails with multimedia information. 








